\section{Introduction}
\label{sec:intro}

Cyberphysical systems (\CPSs) have become the connective tissue of
modern society. Electric power grids, water and wastewater utilities,
oil and gas pipelines, manufacturing plants, hospitals, building
automation systems, and transportation networks all rely on software
that observes and acts upon the physical world. Over the last two
decades, these systems have become increasingly networked, increasingly
dependent on commodity computing components, and increasingly reachable
from the same internet that carries commodity cybercrime. At the same
time, the consequences of a successful attack against them have grown
correspondingly severe: a compromise of a chemical plant controller,
a medical infusion pump, or a substation protection relay is not a
data breach but a safety incident. Recent years have made this concrete.
Ransomware has shut down fuel pipelines and disrupted hospitals for
weeks; attackers have altered the chemistry of municipal water supplies;
and nation-state adversaries have been observed pre-positioning inside
U.S.\ critical-infrastructure networks. The defensive side of this
arms race has not kept pace.

\subsection{The Problem: Why CPS Security Assessment is Hard}
\label{sec:intro-problem}

Assessing the security posture of a cyberphysical system is categorically
harder than assessing a typical enterprise IT environment, for four
reinforcing reasons.

\paragraph{Heterogeneity and non-standard components.}
A single \cps deployment typically contains a mixture of commodity
servers, programmable logic controllers (PLCs), remote terminal units
(RTUs), human-machine interfaces (HMIs), specialized sensors and
actuators, industrial network equipment (such as serial-to-Ethernet
gateways), and safety systems, each speaking its own protocol (Modbus,
DNP3, OPC~UA, IEC~61850, EtherCAT, BACnet, and dozens of proprietary
variants). Many of these components run firmware for which no source
code, documentation, or test environment is available, and standard
security tooling built for x86 Linux simply does not apply.

\paragraph{Operational sensitivity.}
The techniques that security teams routinely apply in enterprise
environments---active network scanning, credential spraying, fuzzing,
exploit validation---cannot safely be used against a running \cps. A
broadcast probe sent to the wrong subnet can reset a controller; a
malformed Modbus request can halt a production line; a short disruption
of a clinical device can injure a patient. This is not a hypothetical
concern: scanning-induced outages have been documented repeatedly in
industrial environments, and many operators now contractually
\emph{prohibit} active assessment of production networks. The result
is that the tests most likely to find real vulnerabilities are the ones
least likely to be permitted.

\paragraph{Dynamism.}
\CPSs are not static. Firmware is patched (or not patched), equipment is
added, removed, or relocated; vendor connections come and go; operating
modes change seasonally or in response to demand; emergency
reconfigurations leave behind temporary-but-permanent bypasses. An
assessment performed at time $t$ can be partially invalidated by time
$t + \Delta$, where $\Delta$ may be days, not years. Yet current
assessment cadences are typically annual, and sometimes much longer.

\paragraph{Scarce expertise.}
The specialists who can credibly evaluate such a system---engineers who
understand both the physical process and the security-relevant software
stack---are few, expensive, and concentrated in a small number of
consultancies and national laboratories. No plausible growth in the
workforce will close this gap at the speed at which the attack surface
is growing.

\subsection{The State of the Art, and Why It Is Insufficient}
\label{sec:intro-sota}

Current practice for \cps security assessment reflects these
constraints. At the high end, a site will undergo a periodic manual
penetration test: a small team of experts spends days or weeks
interviewing operators, reading documentation, passively sniffing
traffic, and cautiously interacting with a limited subset of the system,
typically producing a report that is consumed once and then shelved. At
the lower end, passive asset inventory tools and industrial IDS
products observe network traffic and raise alerts against known-bad
patterns, but they do not reason about how a sophisticated adversary
might combine individually-innocuous weaknesses into an end-to-end
breach. In between lie isolated-lab testbeds in which a vendor's
equipment is analyzed in depth, but these labs are expensive to build
and inevitably diverge from any particular production deployment.

Three limitations are common to all of these approaches. First, they
are \emph{sporadic}: a site may be analyzed once a year, or once a
decade, while the threat landscape and the site itself evolve
continuously. Second, they are \emph{shallow with respect to
composition}: even excellent individual-component analyses stop short
of the multi-step attacks that characterize real incidents, because
such reasoning is performed (if at all) in the head of a human expert
whose time is limited. Third, they fail to exploit the central asset
of the defender---knowledge of the actual system's configuration,
documentation, and behavior---at scale, because there is no
infrastructure to translate that knowledge into automated adversarial
testing.

Meanwhile, attackers are moving in the opposite direction. Commodity
offensive tooling now includes AI-assisted reconnaissance, vulnerability
triage, and exploitation support. Recent evaluations of frontier
language models demonstrate measurable, rapid progress on multi-step
cyber-attack scenarios that only a few years ago would have required
a skilled human operator. The asymmetry is stark: the defender's
workflow is still artisanal, while the attacker's is increasingly
automated.

\subsection{The Proposed Approach: \sysname}
\label{sec:intro-approach}

We propose \sysname, an agentic framework for the continuous,
non-disruptive, in-depth security assessment of cyberphysical systems.
The central idea is to move the adversarial activity off of the real
system and onto a \emph{faithful, continuously-updated digital twin},
where arbitrary---and otherwise unsafe---analyses can be performed
24~hours a day, 7~days a week, without any risk to the physical process.
Human experts remain in the loop, but their scarce attention is applied
where it matters most: confirming high-confidence findings, resolving
ambiguity that agents cannot, and validating suspected vulnerabilities
on the real system once a twin-based analysis has made the case that
validation is worth attempting.

\sysname is organized around four cooperating teams of AI agents, which
together form a closed feedback loop rather than a linear pipeline:

\begin{itemize}[leftmargin=1.3em,itemsep=0.15em,topsep=0.3em]
  \item \textbf{Knowledge Base (KB) Agents} ingest all forms of
  documentation about the target \cps---configuration files, logs,
  network diagrams, equipment manuals, software bills of materials,
  infrastructure-as-code artifacts, security policies, and
  change-management records---and organize this evidence into a hybrid
  knowledge base combining a graph database (for structured facts such
  as topology and component roles) and a vector database (for
  natural-language content that supports retrieval-augmented agent
  reasoning). KB Agents proactively interact with human operators when
  they detect missing or contradictory information.

  \item \textbf{Observer Agents} use the KB to construct an
  \emph{observation plan} that establishes strictly passive vantage
  points across the \cps: network taps, log readers, container-based
  sensors in specific subnets, service-mesh integrations, and eBPF
  probes where permitted. Observer Agents confirm or refute the KB's
  picture of the system and discover components, software versions,
  and flows that documentation missed. Their observations are
  continuous, so the KB evolves with the real system.

  \item \textbf{Replication Agents} consume the KB and observations to
  synthesize a digital twin: a virtualized, emulated, or containerized
  replica that reproduces the system's components, their software
  configuration, their network topology, and representative data. For
  specialized hardware, Replication Agents build on the PIs' prior
  firmware re-hosting work---Conware and HALucinator---and, where even
  that is not possible, synthesize \emph{mock} components whose behavior
  is derived from observed API and network interactions. Crucially,
  fidelity is an explicit, quantified parameter of the replication
  algorithm rather than an implicit assumption.

  \item \textbf{Red Team Agents} continuously attack the digital twin.
  They compose \emph{component-level} analyses (fuzzing, static
  analysis, symbolic execution, taint tracking) with \emph{system-level}
  multi-step attack planning, building on our ChainReactor work on
  AI-planning-based privilege escalation and on the Artiphishell system
  developed for the DARPA AIxCC competition. Each candidate attack is
  grounded in a precondition/postcondition model expressed in PDDL, and
  its goal is to reach a nominated ``flag'' whose compromise has a
  clear security interpretation.
\end{itemize}

A finding produced by a Red Team Agent is not the end of the process.
Every confirmed attack on the twin is presented to a human security
analyst for triage. If the analyst judges the finding to be a
true-positive candidate, they validate it on the real system using an
appropriately cautious procedure, producing either a confirmed
vulnerability or additional evidence that the twin diverges from
reality. If the finding is judged to be a false positive caused by an
inaccurate twin, the fact of that inaccuracy is fed back to the KB,
Observer, and Replication Agents, which jointly plan updates to
documentation, observations, and configuration. Over time, this loop
drives both the twin's fidelity and the Red Team's precision upward.

\subsection{Contributions and Research Questions}
\label{sec:intro-contribs}

This project will produce the following concrete contributions.

\begin{enumerate}[leftmargin=1.5em,itemsep=0.15em,topsep=0.3em]
  \item A design for a \emph{hybrid, agent-native knowledge base} for
  cyberphysical systems that combines structured graph representations
  of topology and component roles with vectorized natural-language
  representations of documentation, policies, and observations, and
  that is resilient to partial, stale, and contradictory inputs.

  \item A suite of \emph{strictly-passive observation techniques} for
  \CPSs---network, log, API, configuration, and security-tool-based---
  and a principled method for synthesizing them into a continuously
  evolving system model.

  \item A \emph{controlled-fidelity digital twin synthesis algorithm}
  in which diversity of component classes, topological patterns, and
  data flows are explicit parameters, enabling minimally representative
  twins of large-scale systems.

  \item An \emph{agentic red-team workflow} that composes component-level
  vulnerability analyses into multi-step attack plans via PDDL-based
  planning, adapted from our ChainReactor and Artiphishell work and
  extended with the emergent reasoning capabilities of modern LLMs.

  \item An \emph{open-source reference implementation} of \sysname,
  together with a set of reference \cps testbeds and benchmarks that
  the research community can use to compare future approaches.
\end{enumerate}

These contributions will answer the following research questions:
\textbf{(RQ1)} Can a team of AI agents construct and maintain a useful,
living model of a cyberphysical system from the evidence that a typical
operator already has?
\textbf{(RQ2)} What is the minimum fidelity a digital twin needs in
order to surface the same classes of vulnerabilities as the real
system, and how can that fidelity be measured?
\textbf{(RQ3)} Can AI-agent red teams, operating against a twin,
discover multi-step vulnerabilities with precision and recall
competitive with---and eventually exceeding---human red teams?
\textbf{(RQ4)} How should human experts be positioned in the loop so
that their scarce attention produces maximum marginal value?

\subsection{Organization of the Proposal}
The remainder of this proposal is organized as follows.
Section~\ref{sec:approach} describes the overall technical approach and
the four agent teams.
Section~\ref{sec:thrusts} details the three research thrusts that
constitute the bulk of the proposed work.
Section~\ref{sec:evaluation} presents our evaluation plan;
Section~\ref{sec:timeline} the project timeline;
Section~\ref{sec:broader} our broader impacts;
Section~\ref{sec:pis} the qualifications of the PIs;
and Section~\ref{sec:prior} results from prior NSF support.
